\section{Multiplicity from Suzuki's inverse compressions}
\label{sec:suzuki-probes}
Suzuki's norm convergence supplies a second finite procedure for recovering
each fixed multiplicity, through the mass of a spectral projection.
His Theorem~3.3 identifies the local inverse with a compression of $C$;
equation~(4.1) proves convergence in Hilbert--Schmidt norm, and
Remark~4.1 gives resolvent convergence away from the ordinates
\cite[pp.~10--12]{suzuki}. His normalized spectral function
converges to $\Delta_\Gamma$, as Theorem~1.2 states. The following
multiplicity recovery is a consequence of that convergence and the
sampling norm. We credit the convergence mechanism to Suzuki.

Let $V_N\subset H$ be the continuous piecewise-linear functions supported
on $[-N,N]$, with mesh $2^{-N}$, zero endpoint values, and integral zero.
They form nested finite-dimensional spaces. Their union is dense in $H$.
Indeed, the fixed-window form is continuous in the $H^1$ norm, so dyadic
interpolation approximates every compact smooth mean-zero function.
Subtract a fixed unit-integral tent times the integral error to keep that
condition exact. General compact smooth functions are approximated in
$q$ by mean-zero ones: for $h\in C_c^\infty$ with $\int h=1$, the
rescaling $h_R(x)=R^{-1}h(x/R)$ satisfies
\[
 q(h_R)\le R^{-2}\|h'\|_1^2\sum_\gamma m_\gamma/\gamma^2\longrightarrow0,
\]
so $f-(\int f)h_R\to f$. Let $P_N$ be the $q$-orthogonal projection
onto $V_N$. Consequently
\begin{equation}
 C_N=P_NCP_N\longrightarrow C\quad\hbox{in }\mathcal S_2.
 \label{eq:finite-compression-limit}
\end{equation}
This is Suzuki's compact-compression argument applied to the indicated
finite spaces.

For a basis $b_1,\ldots,b_d$, all finite data are Weil pairings:
\[
 G=(q(b_j,b_i))_{i,j},\quad
 H_N=(q(\mathcal Pb_j,b_i))_{i,j},\quad
 c=(q(\phi,b_i))_i.
\]
Here $\mathcal P$ is the compact primitive defined above. The matrix
$S_N=G^{-1/2}H_NG^{-1/2}$ represents $C_N$, and $G^{-1/2}c$ represents
$P_N\phi$. The explicit formula uses only prime powers with logarithm
at most twice the largest support radius. Positivity of $G$ permits
certified computation of its inverse square root, by increasing precision
until a positive lower eigenvalue enclosure is obtained. Each finite
search terminates; no uniform condition number is claimed.

\begin{remark}[Fixed multiplicity recovery from Suzuki's convergence]
\label{rem:suzuki-multiplicity}
Fix a compact probe $\phi$ with $\widehat\phi(t)\ne0$ for real $t$.
For a fixed ordinate $\gamma$, let $E_{N,\gamma}$ be the finite spectral
projection converging to the isolated eigenvalue $1/\gamma$ of $C$, and
let $\theta_{N,\gamma}\to1/\gamma$ be its eigenvalue. Then
\[
 \frac{\|E_{N,\gamma}P_N\phi\|_q^2}
      {|\widehat\phi(1/\theta_{N,\gamma})|^2}
           \longrightarrow m_\gamma.
\]
Nearest-integer rounding therefore recovers $m_\gamma$ for every
sufficiently large $N$, with the threshold allowed to depend on $\gamma$.
\end{remark}

\begin{proof}
In the sampling realization, the unit coordinate vector at $\gamma$ has
value $m_\gamma^{-1/2}$ at that coordinate. Thus the squared mass of
$\phi$ on its eigenspace is
$m_\gamma|\widehat\phi(\gamma)|^2$. Norm convergence of $C_N$ gives
norm convergence of the spectral projection by the resolvent identity
on a small circle about $1/\gamma$. The limiting projection has rank
one; projections within distance less than one have equal finite rank,
so the corresponding finite eigenvalue is eventually simple.
Moreover $E_{N,\gamma}P_N\phi=E_{N,\gamma}\phi$.
The numerator therefore converges to the displayed coordinate mass,
and continuity and nonvanishing of the denominator prove the assertion.
A convergent sequence with integer limit eventually lies less than
$1/2$ from that integer.
\end{proof}

We now quantify the two errors in this division. The probe
$\phi=h*\widetilde h$, where
$h(x)=e^{-x}{\bf1}_{[0,1]}(x)$ and $\widetilde h(x)=h(-x)$, is real,
even, supported on $[-1,1]$, and belongs locally to $H^1_0$.
Writing $r=e^{-1}$, its transform is exactly
\[
 p(t)=\widehat\phi(t)
 =\left|\frac{1-e^{-1-it}}{1+it}\right|^2
 =\frac{1+r^2-2r\cos t}{1+t^2}.
\]
In particular
\begin{equation}
 p(t)\ge\frac{(1-r)^2}{1+t^2},\qquad
 |(\log p)'(t)|\le C_0:=1+\frac{2r}{(1-r)^2}.
 \label{eq:probe-bounds}
\end{equation}
The first inequality follows from $|1-re^{-it}|\ge1-r$. For the
second, differentiate the numerator and denominator; their logarithmic
derivatives have absolute values at most $2r/(1-r)^2$ and
$2|t|/(1+t^2)\le1$, respectively.

\begin{proposition}[A certified spectral rounding bound]
\label{prop:spectral-rounding}
Let $\lambda=1/\gamma$ and
$d=\operatorname{dist}(\lambda,\sigma(C)\setminus\{\lambda\})$,
where zero is included in the spectrum. Let $\widetilde C_N$ be a
selfadjoint finite approximation supported on $V_N$, with
$\|\widetilde C_N-C\|\le\varepsilon<d/4$. Let
$\widetilde a_N$ approximate $P_N\phi$ with error at most $\eta$,
and put $B=\|\phi\|_q$.
The spectral point $\widetilde\lambda$ enclosed by
$|z-\lambda|=d/2$ is simple, and its projection $\widetilde E$ satisfies
\[
 \|\widetilde E-E\|\le
 \delta:=\frac{2\varepsilon}{d-2\varepsilon},\qquad
 |\widetilde\gamma-\gamma|\le
 \xi:=\frac{\varepsilon\gamma^2}{1-\varepsilon|\gamma|},
 \quad\widetilde\gamma=1/\widetilde\lambda.
\]
The computed mass $u=\|\widetilde E\widetilde a_N\|_q^2$ satisfies
\[
 |u-m_\gamma p(\gamma)^2|
       \le\delta B^2+2B\eta+\eta^2=:\mathcal E.
\]
For $|\gamma|\le T$, the resulting multiplicity division has error
\begin{equation}
 \left|\frac{u}{p(\widetilde\gamma)^2}-m_\gamma\right|
 \le
 \frac{(1+(T+\xi)^2)^2}{(1-r)^4}\mathcal E
       +m_\gamma\bigl(e^{2C_0\xi}-1\bigr).
 \label{eq:spectral-rounding-error}
\end{equation}
Adding the final certified arithmetic error gives a rounding certificate
as soon as the sum is strictly less than $1/2$.
For a certificate the unknown $m_\gamma$ on the right may be replaced
by any certified upper bound. One always has the probe bound
$m_\gamma\le B^2(1+T^2)^2/(1-r)^4$; a bound for the total zero count
in a certified packet is another choice.
\end{proposition}

\begin{proof}
On the circle of radius $d/2$, the resolvent of $C$ has norm at most
$2/d$. The Neumann series bounds that of $\widetilde C_N$ by
$1/(d/2-\varepsilon)$. Integrate their resolvent identity around the
circle; its length is $\pi d$, giving the bound $\delta$.
Since $\varepsilon<d/4$, this bound is less than one. Both contour
projections are orthogonal, so their ranges have the same dimension,
namely one. The eigenvalue perturbation bound
$|\widetilde\lambda-\lambda|\le\varepsilon$ follows from the
selfadjoint spectral distance estimate, or the min--max principle.
Since $d\le|\lambda|$ and $\varepsilon<d/4$, taking reciprocals is
legitimate and gives the bound $\xi$.

Because $\widetilde E$ is supported on $V_N$,
$\widetilde EP_N\phi=\widetilde E\phi$. The difference of the two
exact probe masses is at most $\delta B^2$. Replacing $P_N\phi$ by
$\widetilde a_N$ contributes at most $2B\eta+\eta^2$ by expanding
the squared norm. This proves the mass estimate without requiring a
separate convergence modulus for $P_N\phi$.
Divide that estimate by $p(\widetilde\gamma)^2$ and use
\eqref{eq:probe-bounds}. Finally, integrating the logarithmic derivative
bound gives
$|p(\gamma)^2/p(\widetilde\gamma)^2-1|
\le e^{2C_0|\gamma-\widetilde\gamma|}-1$.
The claimed division error follows.
\end{proof}

\begin{proposition}[A local residual certificate for the probe]
\label{prop:local-probe-residual}
Let $v\in V_N$ have $q(v)=1$, and put
\[
 \mu=\langle Cv,v\rangle_q,\qquad
 \rho=\|(C-\mu)v\|_q.
\]
Assume $|\mu|>\rho$, and let $\lambda$ be the unique spectral point
in $[\mu-\rho,\mu+\rho]$.  Let $d>0$ be a lower bound for its
distance to the rest of $\sigma(C)$, including zero, with
$\rho<d/2$.  Define
\[
 \delta_v=\frac{\rho}{d-\rho},\qquad
 \omega=\min\left\{\rho,
             \frac{\rho^2(d-\rho)}{d(d-2\rho)}\right\},\qquad
 \xi_v=\frac{\omega}{|\mu|(|\mu|-\omega)}.
\]
Then, writing $E_\lambda$ for the spectral projection and
$\widetilde\gamma=1/\mu$, $\gamma=1/\lambda$, one has
\[
 \|v\otimes v-E_\lambda\|\le\delta_v,
 \qquad |\mu-\lambda|\le\omega,
 \qquad |\widetilde\gamma-\gamma|\le\xi_v.
\]
The finite probe mass $u=|\langle\phi,v\rangle_q|^2$ therefore obeys
\begin{equation}
 \left|\frac{u}{p(\widetilde\gamma)^2}-m_\gamma\right|
 \le\frac{(1+(T+\xi_v)^2)^2}{(1-r)^4}\delta_v B^2
             +m_\gamma(e^{2C_0\xi_v}-1)
 \quad (|\gamma|\le T),
 \label{eq:local-probe-rounding}
\end{equation}
with $r,B,C_0$ as in \eqref{eq:probe-bounds} and
Proposition~\ref{prop:spectral-rounding}.  The same upper bounds for
$m_\gamma$ and the same strict $1/2$ rounding threshold apply.
\end{proposition}

\begin{proof}
The spectral theorem gives
$\operatorname{dist}(\mu,\sigma(C))\le\rho$.  Since $|\mu|>\rho$,
the enclosed spectral point is nonzero.  Write
$v=\sum_j a_je_j$ in the orthonormal coordinate basis of $C$, and
let $w=|a_\lambda|^2$.  Every other spectral point obeys
$|\lambda_j-\mu|\ge d-\rho$.  Consequently
\[
 1-w\le\frac{\rho^2}{(d-\rho)^2}.
\]
The distance between the two rank-one orthogonal projections is
$\sqrt{1-w}$, which gives the first inequality.  The Rayleigh
identity yields
\[
 (\mu-\lambda)w
      =\sum_{j\ne\lambda}(\lambda_j-\mu)|a_j|^2.
\]
Cauchy--Schwarz bounds the absolute value on the right by
$\rho\sqrt{1-w}\le\rho^2/(d-\rho)$.  Dividing by
$w\ge1-\rho^2/(d-\rho)^2>0$ proves the second expression in the
definition of $\omega$; spectral inclusion already gives the first.
Taking reciprocals proves the bound $\xi_v$.

The projection estimate gives
$|u-\|E_\lambda\phi\|_q^2|\le\delta_v B^2$.
In the sampling space that exact mass is $m_\gamma p(\gamma)^2$.
Division and the logarithmic derivative bound for $p$ give
\eqref{eq:local-probe-rounding} exactly as in the preceding proof.
\end{proof}

All data in this residual can be computed from finite Weil pairings.
In addition to $G,H_N,c$, form the primitive Gram matrix
\[
 J_N=(q(\mathcal Pb_j,\mathcal Pb_i))_{i,j}.
\]
For a nonzero coefficient vector $x$, take
$v=\sum_jx_jb_j/(x^*Gx)^{1/2}$.  Then the exact formulas are
\begin{equation}
 \mu=\frac{x^*H_Nx}{x^*Gx},\qquad
 \rho^2=\frac{x^*J_Nx}{x^*Gx}-\mu^2,\qquad
 u=\frac{|x^*c|^2}{x^*Gx}.
 \label{eq:residual-arithmetic}
\end{equation}
Thus no separate approximation to $P_N\phi$ is needed.  The primitive
of a compact function with integral zero has the same compact support,
so the explicit formula for $J_N$ has a finite prime-power cutoff too.
Certified interval errors in \eqref{eq:residual-arithmetic} and in the
final divided scalar are included in the rounding inequality.

For translated tents this new Gram matrix has a direct scalar recipe.
Let $t$ be the tent centered at zero and let
$b_i(x)=t(x-ih)-t(x)$.  Define the even function
\[
 A_2(d)=\sum_{\gamma\in\Gamma}
      m_\gamma|\widehat t(\gamma)|^2
                      \frac{1-\cos(\gamma d)}{\gamma^2}.
\]
Expanding the two translation factors gives
\[
 (J_N)_{ij}=A_2(ih)+A_2(jh)-A_2((i-j)h).
\]
The inverse Fourier transform of
$(1-\cos(\gamma d))/\gamma^2$ is
$\tfrac12(|d|-|x|)_+$.  Therefore the geometric Weil functional
computes $A_2(d)$ from the compact test
$\tfrac12(|d|-|x|)_+*(t*\widetilde t)$.
It is piecewise polynomial of degree five, with explicitly known
breakpoints.  This identifies the additional finite integrals required
by a residual certificate.

For comparison with the original piecewise-polynomial probe, set
$\phi_0(x)=(1-|x|)_++(\sqrt2-|x|)_+$. Its transform obeys the
following fully explicit bounds:
\[
 p_0(t)=\frac{2(1-\cos t)+2(1-\cos(\sqrt2t))}{t^2},\quad p_0(0)=3,
\]
\begin{equation}
 p_0(t)\ge
 \begin{cases}
 12/\pi^2,&|t|\le\pi/\sqrt2,\\[2pt]
 \displaystyle\frac{16}{3t^2
     (\sqrt2|t|/\pi+(\sqrt2+1)/2)^2},&|t|>\pi/\sqrt2.
 \end{cases}
 \label{eq:two-tent-bound}
\end{equation}
Here is the arithmetic proof. By evenness take $t>0$, and choose
integers $n,m$ nearest to $t/(2\pi),\sqrt2t/(2\pi)$, respectively.
Put $e_1=t/(2\pi)-n$, $e_2=\sqrt2t/(2\pi)-m$. Above the stated
threshold these integers are not both zero. The nonzero integer
$m^2-2n^2$ has absolute value at least one, while
$|m+\sqrt2n|\le D:=\sqrt2t/\pi+(\sqrt2+1)/2$. Therefore
$|\sqrt2e_1-e_2|=|m-\sqrt2n|\ge1/D$, and
$e_1^2+e_2^2\ge1/(3D^2)$. Concavity gives
$\sin(\pi|e_j|)\ge2|e_j|$, since $|e_j|\le1/2$. Using
$1-\cos(2\pi e_j)=2\sin^2(\pi e_j)$ proves the second bound.
For $0<t\le\pi/\sqrt2$, apply the same sine inequality directly
to $t/2$ and $\sqrt2t/2$; it gives the first bound.
The value at zero follows by continuity. Thus division by $p_0^2$
has an explicit loss of order $T^8$, while the exponential
self-correlation probe has the order $T^4$ bound in
\eqref{eq:spectral-rounding-error}.

\subsection{A common bounded height and the remaining quantifiers}
This comparison identifies the separation data that both contour and
spectral certification must control. Let $T$ exceed the first ordinate,
and let $s(2T)$ be a certified positive lower bound for the distances
between distinct positive ordinates at most $2T$. For $0<\gamma\le T$,
\begin{equation}
 \operatorname{dist}\bigl(\gamma^{-1},
               \sigma(C)\setminus\{\gamma^{-1}\}\bigr)
 \ge\min\{s(2T)/(2T^2),\,1/(2T)\}.
 \label{eq:reciprocal-separation}
\end{equation}
For another ordinate $0<\eta\le2T$, the identity
$|\gamma^{-1}-\eta^{-1}|=|\gamma-\eta|/(\gamma\eta)$ proves the
first bound. Ordinates above $2T$, negative ordinates, and the spectral
point zero have distance at least $1/(2T)$.
This is the exact conversion from an ordinate gap to the spectral gap.
The same certified ordinate gaps permit disjoint disks of radius less
than one third of the gap in the contour calculation.

The interval selection in Proposition~\ref{prop:local-probe-residual}
can also use these gaps directly.  Write $d_T$ for the right side of
\eqref{eq:reciprocal-separation}.  A positive Rayleigh value with
$\mu-\rho\ge1/T$ and $2\rho<d_T$ has a spectral point in its
residual interval, by the spectral theorem.  That point corresponds to
an ordinate at most $T$, and the gap bound makes it the unique point
in the interval.  Thus the local residual test needs no additional
unknown spectral projection.  For example, an ordinate separation
exceeding $6\cdot10^{-10}$ through height $90$ yields
$d_{45}>1.48\cdot10^{-13}$.  Rounding still requires the separately
certified primitive Gram residual and the divided mass error in
\eqref{eq:local-probe-rounding}; separation alone does not bound them.

The common range consists of the heights for which the contour's
boundary and Taylor errors are certified and the right side of
\eqref{eq:spectral-rounding-error} or
\eqref{eq:local-probe-rounding}, uniformly over the enclosed
ordinates, is less than $1/2$. Suzuki's convergence supplies eventual
correctness at each fixed ordinate, and hence simultaneously below
each fixed height, because there are only finitely many such ordinates.
The stopping index can depend on the height; this convergence statement
does not specify a certified stopping index or assert that the recovered
integers equal one. An all-height certificate of simplicity must supply
an unbounded sequence of heights, finite packets covering all ordinates
below each height, certified errors less than $1/2$, and recovered
integers all equal to one. The contour boundary, Taylor and
leading-coefficient bounds, or the spectral residual and separation
bounds, supply the corresponding finite error tests. A finite
successful list recovers the multiplicities included in that list.

One exact identity also measures the unresolved global mass. The
Hilbert--Schmidt orthogonal projection $C\mapsto P_NCP_N$ gives
\[
 \|C-C_N\|_{\mathcal S_2}^2
      =\|C\|_{\mathcal S_2}^2-\operatorname{tr}(C_N^2).
\]
Here and below $\Gamma$ includes both signs of every distinct ordinate.
On the other hand the Hadamard product gives
$-\Xi''(0)/\Xi(0)=\sum_\gamma m_\gamma/\gamma^2$.
Consequently the computable decreasing quantities
\[
 \mathcal D_N=-\frac{\Xi''(0)}{\Xi(0)}
                      -\operatorname{tr}(C_N^2)
\]
satisfy the exact formula
\begin{equation}
 \mathcal D_N=
 \|C-C_N\|_{\mathcal S_2}^2+
           \sum_{\gamma\in\Gamma}\frac{m_\gamma-1}{\gamma^2}.
 \label{eq:trace-multiplicity-defect}
\end{equation}
The sum in \eqref{eq:trace-multiplicity-defect} is over signed ordinates,
and equals $2\sum_{\gamma>0}(m_\gamma-1)/\gamma^2$.
Their limit is the inverse-square weighted multiplicity excess.
Thus $\mathcal D_N\to0$ is another equivalent formulation of simplicity;
the identity proves the equivalence and does not establish that limit.
To verify the orthogonal-projection identity, choose an orthonormal
basis beginning with one for $V_N$: the squared Hilbert--Schmidt norm
is the sum of squared matrix entries, and compression retains exactly
the entries in the finite square. Increasing $N$ increases this square
norm. The weighted sum for $\Xi''(0)$ is obtained by differentiating
the paired product at zero, where absolute convergence is supplied by
$\sum m_\gamma/\gamma^2<\infty$.

\begin{proposition}[A trace certificate for simplicity below a height]
\label{prop:trace-height}
Let $\mathcal D_N$ be the preceding trace defect and let
$E(T)=\sum_{0<\gamma\le T}(m_\gamma-1)$, with the sum over distinct
positive ordinates. For every $N$ and $T>0$,
\begin{equation}
 E(T)\le\frac{T^2}{2}\mathcal D_N.
 \label{eq:trace-height}
\end{equation}
In particular, a certified upper bound
$\mathcal D_N\le U<2/T^2$ certifies simplicity of every zero with
$0<\gamma\le T$. When $U>0$, it certifies every height
$T<\sqrt{2/U}$; $U=0$ certifies simplicity at every height.
The following conditions are equivalent:
\[
 \text{all nontrivial zeros are simple},\qquad
 \mathcal D_N\longrightarrow0,\qquad
 \text{for every }T>0\text{ some }N\text{ has }\mathcal D_N<2/T^2.
\]
\end{proposition}

\begin{proof}
Every summand of the multiplicity excess and the squared norm in
\eqref{eq:trace-multiplicity-defect} is nonnegative. Restrict its signed
sum to $0<|\gamma|\le T$ and pair the two signs. Since
$\gamma^{-2}\ge T^{-2}$ there,
\[
 \mathcal D_N\ge2\sum_{0<\gamma\le T}
       \frac{m_\gamma-1}{\gamma^2}\ge\frac{2E(T)}{T^2}.
\]
The quantity $E(T)$ is a nonnegative integer, so an upper bound strictly
less than one forces it to vanish. This proves the certificate.
The decreasing limit of $\mathcal D_N$ is the weighted excess in
\eqref{eq:trace-multiplicity-defect}; it vanishes precisely when every
nonnegative integer $m_\gamma-1$ vanishes. This proves the first
equivalence. A sequence tending to zero eventually falls below each
positive threshold $2/T^2$. Conversely, existence of an $N$ below each
such threshold, together with monotonicity, makes the limit zero.
\end{proof}

The finite certificate uses the compression matrices and the constant
$S=-\Xi''(0)/\Xi(0)$. For instance, certified bounds
$S\le S^+$ and $\operatorname{tr}(C_N^2)\ge L_N$ give the usable upper
bound $U=S^+-L_N$. It needs neither an isolating circle for each
reciprocal eigenvalue nor division by an individual probe value.
The trace certificate supplies an upper bound for the aggregate
multiplicity excess. The local contour minors give the multiplicities
inside a chosen contour through their successive vanishing orders,
and the probe gives the multiplicity at an isolated spectral coordinate.
The exact formulas above relate these outputs to the same local
integers $m_\gamma$. A multiple zero above the requested height can
keep the trace defect large, so failure of its finite inequality says
nothing by itself about simplicity below that height. No decay estimate
for $\mathcal D_N$ is established here.
